
\documentclass[11pt,a4paper]{article}
\usepackage[utf8]{inputenc}
\usepackage[T1]{fontenc}
\usepackage[spanish]{babel}
\usepackage{geometry}
\usepackage{graphicx}
\usepackage{array}
\usepackage{tabularx}
\usepackage{makecell}
\usepackage{hyperref}
\usepackage{xcolor}
\usepackage{fancyhdr}
\usepackage{titlesec}
\usepackage{helvet}
\usepackage{ragged2e}
\usepackage{enumitem}
\usepackage{setspace}
\usepackage{amsmath}
\usepackage{booktabs}
\usepackage{float}
\usepackage{listings}
\usepackage{hyperref}
\usepackage{xurl}
\usepackage{microtype}
\usepackage{float}
\emergencystretch=2em
\renewcommand{\familydefault}{\sfdefault}
\geometry{a4paper,left=3cm,right=3cm,top=2.35cm,bottom=2.35cm,headheight=1.9cm,headsep=0.45cm}
\setlength{\parindent}{0pt}
\setlength{\parskip}{0.3em}
\setlength{\tabcolsep}{5pt}
\renewcommand{\arraystretch}{1.13}
\pagestyle{fancy}
\fancyhf{}
\fancyhead[L]{\IfFileExists{reporte_tecnico_overleaf_logo.png}{\includegraphics[height=1.35cm]{reporte_tecnico_overleaf_logo.png}}{\small Universidad Nacional de Loja}}
\fancyhead[R]{\small FEIRNNR -- Carrera de Computación}
\fancyfoot[C]{\thepage}
\renewcommand{\headrulewidth}{0.5pt}
\renewcommand{\footrulewidth}{0pt}
\titleformat{\section}{\bfseries\fontsize{15}{17}\selectfont}{\thesection.}{0.8em}{}
\titlespacing*{\section}{0pt}{1.0ex}{0.3ex}
\newcolumntype{L}[1]{>{\RaggedRight\arraybackslash}p{#1}}
\newcolumntype{Y}{>{\RaggedRight\arraybackslash}X}
\lstset{basicstyle=\ttfamily\footnotesize,breaklines=true,frame=single,showstringspaces=false,columns=fullflexible}
\hypersetup{hidelinks}

\usepackage{calc}
\usepackage{fvextra}
\providecommand{\tightlist}{\setlength{\itemsep}{0pt}\setlength{\parskip}{0pt}}
\begin{document}
\thispagestyle{fancy}
\begin{center}
{\bfseries\fontsize{21}{24}\selectfont Reporte Técnico de Actividades\\ Práctico-Experimentales Nro. 004\par}
\vspace{0.25cm}
{\large Implementación de seguridad, validaciones y pruebas de API}
\end{center}

\section{Datos de Identificación del Estudiante y la Práctica}
\begin{tabularx}{\textwidth}{|L{0.43\textwidth}|Y|}
\hline
\textbf{Nombre del estudiante(s)} & [PENDIENTE: integrantes reales] \\ \hline
\textbf{Asignatura} & Desarrollo Basado en Plataformas \\ \hline
\textbf{Ciclo} & 5 A \\ \hline
\textbf{Unidad} & 1 \\ \hline
\textbf{Resultado de aprendizaje} & Discute cómo los estándares Web influyen en el desarrollo de software, bajo los principios de solidaridad, transparencia, responsabilidad y honestidad. \\ \hline
\textbf{Práctica Nro.} & 004 \\ \hline
\textbf{Título de la Práctica} & Implementación de seguridad, validaciones y pruebas de API \\ \hline
\textbf{Nombre del Docente} & Edison Leonardo Coronel Romero \\ \hline
\textbf{Fecha} & Semana 4; [PENDIENTE: fecha real] \\ \hline
\textbf{Horario} & Por definir según horario académico \\ \hline
\textbf{Lugar} & Aula / laboratorio asignado \\ \hline
\textbf{Tiempo planificado en el Sílabo} & 3 horas \\ \hline
\end{tabularx}
\textbf{Estado:} preparación documental y artefactos; ejecución funcional, corrección y despliegue pendientes.\par

\section{Objetivo(s) de la
Práctica}\label{APE4-semana-4-objetivos-de-la-pruxe1ctica}

\begin{itemize}
\tightlist
\item
  Implementar mecanismos de autenticación y autorización coherentes con
  los actores, roles y recursos del proyecto.
\item
  Aplicar validaciones de entrada, reglas de negocio y manejo seguro de
  errores.
\item
  Verificar la protección de datos y secretos mediante prácticas básicas
  de seguridad.
\item
  Implementar y ejecutar pruebas reproducibles de API, integración y
  controles de acceso.
\item
  Corregir hallazgos identificados durante la revisión de seguridad y
  mantener trazabilidad con requisitos y amenazas.
\end{itemize}

La documentación no convierte controles Spring previos en implementación
nueva de DRF. Los objetivos de ejecución/corrección funcional no se dan
por cumplidos.

\section{Materiales, Reactivos, Equipos y
Herramientas}\label{APE4-semana-4-materiales-reactivos-equipos-y-herramientas}

Insumos: modelo/matriz ACD4 y fuentes compartidas, requisitos/AC AA2,
OpenAPI y evidencia APE3, pruebas existentes de identidad/catálogo.
Stack real: Java/Spring/PostgreSQL, SQLite auxiliar anterior; la guía
exige Python/Django/DRF y pytest-django, no implementados/ejecutados
aquí. Se prepara un cliente pytest, que no es un backend Django.

Variables de entorno son mecanismo previsto de secretos; placeholders no
son valores reales. Git/GitHub existen, GitFlow y contribución están
pendientes. Actions del repositorio usa workflow\_dispatch, no se activa
desde esta entrega. No hay Docker/Podman ni un sandbox conforme
acreditado; WSL disponible no basta para afirmarlo. No se instalan
dependencias ni se ejecuta target fuera del aislamiento exigido.

\section{Procedimiento / Metodología
Ejecutada}\label{APE4-semana-4-procedimiento--metodologuxeda-ejecutada}

Inicio: se seleccionan sesión, catálogo y preparación de proceso; se
relacionan actores/capacidades, T-01--09, SEC-01--05 y AC existentes.
ACD4 no tiene aún retroalimentación humana real. La revisión es estática
y usa evidencia histórica, no una nueva campaña local sobre servicio
compartido.

\textbf{Paso }: 1 Autenticación

\textbf{Trabajo efectivamente preparado/revisado }: SessionService,
BCrypt, hashes, cookie, expiración/revocación y ventana; VT-01--03

\textbf{Estado de cumplimiento }: Implementado previamente en Spring;
nueva ejecución y DRF pendientes.

\medskip\hrule\medskip

\textbf{Paso }: 2 Autorización

\textbf{Trabajo efectivamente preparado/revisado
}: ScopeAuthorizationService, CAREER\_ADMIN/READ, PROCESS\_MANAGE;
fixtures lector A/carrera B

\textbf{Estado de cumplimiento }: Fuente y matriz; prueba funcional
nueva pendiente.

\medskip\hrule\medskip

\textbf{Paso }: 3 Validaciones

\textbf{Trabajo efectivamente preparado/revisado }: DTO, sintaxis, UUID,
fechas, relaciones, duplicado y ETag; VT-07

\textbf{Estado de cumplimiento }: Controles previos inspeccionados, sin
cambio de negocio.

\medskip\hrule\medskip

\textbf{Paso }: 4 Datos/secretos

\textbf{Trabajo efectivamente preparado/revisado }: Escaneo acotado de
texto y nombres históricos, env/prod/.gitignore; reporte sin valores

\textbf{Estado de cumplimiento }: Comprobación documental limitada, no
certificación de ausencia total.

\medskip\hrule\medskip

\textbf{Paso }: 5 Errores

\textbf{Trabajo efectivamente preparado/revisado }: Problem Details,
estados y correlación; desafío401/sobre pendientes

\textbf{Estado de cumplimiento }: No se corrige protocolo ni se migra
contrato por documentación.

\medskip\hrule\medskip

\textbf{Paso }: 6 Tests

\textbf{Trabajo efectivamente preparado/revisado }: Siete tests pytest
acotados en \path|pruebas/test_seguridad_api.py|; matriz VT-01--12

\textbf{Estado de cumplimiento }: Código de prueba propuesto; no
ejecutado. Sin gate se omite y no cuenta como aprobado.

\medskip\hrule\medskip

\textbf{Paso }: 7 Verificación segura

\textbf{Trabajo efectivamente preparado/revisado }: Requisitos de
sandbox y datos sintéticos en PLAN\_PRUEBAS

\textbf{Estado de cumplimiento }: No se ejecuta target por falta de
aislamiento completo. No servicios ajenos.

\medskip\hrule\medskip

\textbf{Paso }: 8 Correcciones

\textbf{Trabajo efectivamente preparado/revisado }: Acciones SEC
concretas y CF-04-01 reservado

\textbf{Estado de cumplimiento }: Sin parche/regresión ejecutada; no hay
antes\ensuremath{\rightarrow}después de seguridad acreditado.

\medskip\hrule\medskip

\textbf{Paso }: 9 Trazabilidad

\textbf{Trabajo efectivamente preparado/revisado
}: Amenaza\ensuremath{\rightarrow}SEC\ensuremath{\rightarrow}RF/RNF\ensuremath{\rightarrow}endpoint\ensuremath{\rightarrow}control\ensuremath{\rightarrow}VT\ensuremath{\rightarrow}estado

\textbf{Estado de cumplimiento }: Matriz elaborada, código/contrato de
producto sin modificar.

\medskip\hrule\medskip

Los tests propuestos incluyen ausencia de credencial, desafío401,
identidad desconocida sintética, falta de CSRF, campos obligatorios,
objeto ajeno y logout. El escenario 401 con desafío debe servir de
regresión cuando se acuerde un mecanismo compatible; no se fuerza JWT.
El plan limita red/host/tamaño/respuestas y datos, pero ese gate no es
la sandbox OS.

Para expiración se selecciona test existente con reloj controlado, no se
espera media hora ni se alteran cuentas ajenas. Para rollback de
auditoría y concurrencia se requiere fixture dedicado; no se extrapolan
las aserciones secuenciales de APE3.

Cierre: falta demostrar flujo protegido al docente, corrección de
seguridad reproducida y contribución versionada real. Un test propuesto
que se sabe que necesita corrección no es un fallo observado nuevo. El
409 de modelo de APE3 corrigió el escenario de representación, no una
vulnerabilidad APE4.

\section{Resultados}\label{APE4-semana-4-resultados}

\textbf{Resultado esperado }: Autenticación/autorización implementadas y
verificadas

\textbf{Resultado de esta entrega }: Código Spring existente
inspeccionado y pruebas fuente localizadas

\textbf{Pendiente }: DRF/decisión tecnológica y ejecución en sandbox.

\medskip\hrule\medskip

\textbf{Resultado esperado }: Validaciones y negocio

\textbf{Resultado de esta entrega }: Matriz sintaxis/negocio y controles
por operación

\textbf{Pendiente }: Ejecutar variantes VT y corregir si fallan.

\medskip\hrule\medskip

\textbf{Resultado esperado }: Secretos fuera del código

\textbf{Resultado de esta entrega }: Configuración externa y escaneo
acotado sin coincidencias de tres patrones

\textbf{Pendiente }: Formatos no cubiertos, historia completa/rotación
si hay indicio.

\medskip\hrule\medskip

\textbf{Resultado esperado }: Errores/API

\textbf{Resultado de esta entrega }: Formato actual y diferencias
trazados

\textbf{Pendiente }: Sobre universal/desafío y pruebas aceptadas.

\medskip\hrule\medskip

\textbf{Resultado esperado }: Pruebas reproducibles

\textbf{Resultado de esta entrega }: Cliente pytest propuesto,
PLAN\_PRUEBAS y VT

\textbf{Pendiente }: Dependencias, sandbox, fixtures,
ejecución/resultados.

\medskip\hrule\medskip

\textbf{Resultado esperado }: Correcciones

\textbf{Resultado de esta entrega }: Acciones/localización/verificación
definidas

\textbf{Pendiente }: Ningún parche funcional nuevo demostrado.

\medskip\hrule\medskip

\textbf{Resultado esperado }: Trazabilidad

\textbf{Resultado de esta entrega }: Matrices compartidas y relación con
AC

\textbf{Pendiente }: Añadir commit/resultados cuando existan.

\medskip\hrule\medskip

\textbf{Resultado esperado }: Contribución individual

\textbf{Resultado de esta entrega }: Nota preparada

\textbf{Pendiente }: Identidad y evidencia personal reales.

\medskip\hrule\medskip

\section{Preguntas de Control}\label{APE4-semana-4-preguntas-de-control}

\subsection{1. ¿Qué diferencia existe entre autenticación y
autorización?}\label{APE4-semana-4-1-quuxe9-diferencia-existe-entre-autenticaciuxf3n-y-autorizaciuxf3n}

Autenticación determina la identidad mediante credencial/sesión;
autorización evalúa capacidad y ámbito sobre una operación/recurso.
SessionService valida token; CareerAdministrationService exige
CAREER\_ADMIN o lectura visible. Una sesión no concede todas las
capacidades.

\subsection{2. ¿Por qué un usuario autenticado no debería poder acceder
automáticamente a todos los
recursos?}\label{APE4-semana-4-2-por-quuxe9-un-usuario-autenticado-no-deberuxeda-poder-acceder-automuxe1ticamente-a-todos-los-recursos}

Porque carreras/procesos están sujetos a ámbito y separación de
responsabilidades. Un lector de A no debe recibir detalle, lista ni
totales de B. UUID y enlace visible no sustituyen la comprobación del
backend.

\subsection{3. ¿Qué diferencia existe entre validación sintáctica y
validación de reglas de
negocio?}\label{APE4-semana-4-3-quuxe9-diferencia-existe-entre-validaciuxf3n-sintuxe1ctica-y-validaciuxf3n-de-reglas-de-negocio}

Sintaxis comprueba formato/tipo/longitud, por ejemplo UUID o fecha.
Negocio verifica unidad activa, período permitido, duplicado normalizado
y transición. Una fecha ISO válida puede exceder la vigencia de la
carrera.

\subsection{4. ¿Qué riesgos existen si se almacenan secretos
directamente en el
repositorio?}\label{APE4-semana-4-4-quuxe9-riesgos-existen-si-se-almacenan-secretos-directamente-en-el-repositorio}

Pueden difundirse en clones, CI e historial aunque se borre el archivo
actual; permitirían uso indebido de cuentas/datos. Usar configuración
externa, mínimo privilegio y rotar/revocar ante exposición confirmada.
No confundir contraseña sintética de fixture con secreto real sin
clasificarla.

\subsection{5. ¿Cómo puede una prueba demostrar que un control de
autorización funciona
correctamente?}\label{APE4-semana-4-5-cuxf3mo-puede-una-prueba-demostrar-que-un-control-de-autorizaciuxf3n-funciona-correctamente}

Comparar mismo recurso/operación con principal autorizado y lector ajeno
de fixture: permitir el primero y rechazar/ocultar el segundo sin efecto
ni datos. Comprobar DB y totales si aplica. Un 401 anónimo no prueba
permiso por objeto ni un mock cubre todo el filtro.

\subsection{6. ¿Qué información no debería exponerse en una respuesta de
error?}\label{APE4-semana-4-6-quuxe9-informaciuxf3n-no-deberuxeda-exponerse-en-una-respuesta-de-error}

Contraseña, token completo, claves, SQL sensible, stack, detalles de
infraestructura o datos ajenos. Devolver código público, detalle seguro
y correlación; registrar diagnóstico protegido. La revisión actual no
prueba todos los logs/500.

\subsection{7. ¿Cómo se relacionan una amenaza, un control de seguridad
y una prueba
verificable?}\label{APE4-semana-4-7-cuxf3mo-se-relacionan-una-amenaza-un-control-de-seguridad-y-una-prueba-verificable}

T-02 acceso ajeno\ensuremath{\rightarrow}permiso por capacidad/scope\ensuremath{\rightarrow}VT-04 con lector A/carrera
B y rechazo sin dato/efecto. Si falla, registrar SEC/CF, cambiar el
punto de control y repetir regresión; conservar resultados y commit
reales, no solo un tick documental.

\section{Conclusiones}\label{APE4-semana-4-conclusiones}

Hay mecanismos Spring fuente y pruebas previas que permiten diseñar una
verificación concreta, no una lista OWASP genérica. Esta entrega no
acredita una nueva implementación DRF, una campaña APE4 ejecutada o un
parche de seguridad. La práctica permanece parcial en esos puntos y en
su cierre docente.

\section{Recomendaciones}\label{APE4-semana-4-recomendaciones}

Resolver tecnología y contrato; establecer sandbox; preparar fixture
PostgreSQL; ejecutar VT con límites y clasificar fallos; aplicar la
corrección mínima en el último punto confiable; reejecutar y versionar
aportes reales. Mantener Actions manual y no usar despliegue ajeno para
estas pruebas.

\section{Bibliografía /
Referencias}\label{APE4-semana-4-bibliografuxeda--referencias}

\href{../soporte-semanas-4-5/FUENTES.md}{Fuentes compartidas}, F-01--12,
OWASP O-01--08, HTTP H-01 y pytest D-05. La guía académica origina los
resultados esperados; los mecanismos fuente no equivalen a seguridad
certificada.

\nocite{*}
\bibliographystyle{IEEEtran}
\bibliography{referencias}
\section{Anexos}\label{APE4-semana-4-anexos}

\url{PLAN_PRUEBAS.md}, prueba pytest propuesta, matrices compartidas,
escaneo de metadatos, registro CF/CRUZADA, nota individual y manifiesto.
Los JSON APE3 originales se referencian, no se recopian con fecha de
Semana4.






\subsection{Fuente propuesta: test\_seguridad\_api.py}
\begin{Verbatim}[breaklines=true,breakanywhere=true,fontsize=\footnotesize]
"""Proposed bounded tests; NOT executed or evidence of DRF implementation.

Run only in the OS-isolated synthetic environment described in PLAN_PRUEBAS.md.
The module skips by default and rejects remote URLs and redirects.
"""
import http.cookiejar
import json
import os
from urllib.error import HTTPError
from urllib.parse import urlsplit
from urllib.request import HTTPRedirectHandler, HTTPCookieProcessor, ProxyHandler, Request, build_opener
import uuid

import pytest

if os.environ.get("SDD_SECURITY_TESTS") != "1":
    pytest.skip("Security runtime checks not explicitly enabled in isolated fixture", allow_module_level=True)

BASE = os.environ.get("SDD_TEST_BASE_URL", "http://127.0.0.1:18094/api/v1").rstrip("/")
url = urlsplit(BASE)
if url.hostname not in ("127.0.0.1", "localhost", "::1") or url.scheme not in ("http", "https"):
    raise RuntimeError("Only an isolated loopback backend is permitted")
if url.username or url.password or url.query or url.fragment:
    raise RuntimeError("Credentials/query/fragment are not allowed in base URL")


class NoRedirects(HTTPRedirectHandler):
    def redirect_request(self, request, fp, code, message, headers, newurl):
        raise RuntimeError("Redirect rejected; remain inside isolated backend")


class Client:
    def __init__(self):
        self.opener = build_opener(ProxyHandler({}), HTTPCookieProcessor(http.cookiejar.CookieJar()), NoRedirects())
        self.csrf = None

    def call(self, method, path, payload=None, csrf=True):
        if not path.startswith("/") or path.startswith("//"):
            raise ValueError("Expected relative API resource")
        headers = {"Accept": "application/json, application/problem+json"}
        if payload is not None:
            headers["Content-Type"] = "application/json"
        if csrf and self.csrf:
            headers["X-CSRF-TOKEN"] = self.csrf
        request = Request(BASE + path, method=method, headers=headers,
                          data=None if payload is None else json.dumps(payload).encode())
        try:
            response = self.opener.open(request, timeout=10)
        except HTTPError as error:
            response = error
        with response:
            raw = response.read(262145)
            if len(raw) > 262144:
                raise RuntimeError("Response exceeds bounded fixture limit")
            return response.status, response.headers, json.loads(raw) if raw else None

    def login(self, prefix):
        user = os.environ.get(prefix + "_USERNAME")
        password = os.environ.get(prefix + "_PASSWORD")
        if not user or not password:
            pytest.skip("Synthetic principal variables are required; no credentials guessed")
        status, _, body = self.call("POST", "/sessions", {"username": user, "password": password})
        assert status == 201, "Synthetic principal login failed"
        self.csrf = body["csrfToken"]


@pytest.fixture
def admin():
    client = Client()
    client.login("SDD_ADMIN")
    try:
        yield client
    finally:
        client.call("DELETE", "/sessions")


def test_vt01_no_credentials_rejected():
    status, _, body = Client().call("GET", "/careers")
    assert status == 401
    assert body["code"] and body["correlationId"]


def test_vt01_challenge_header_regression():
    status, headers, _ = Client().call("GET", "/careers")
    assert status == 401
    assert headers.get("WWW-Authenticate"), "Resolve HTTP-AA-01 using the accepted authentication scheme"


def test_vt02_unknown_synthetic_identity_rejected():
    status, _, body = Client().call("POST", "/sessions", {
        "username": "missing-" + uuid.uuid4().hex[:12], "password": "Synthetic-Invalid-Only-2026"})
    assert status == 401
    assert body["code"] == "AUTHENTICATION_FAILED"


def test_vt06_missing_csrf_rejected(admin):
    status, _, _ = admin.call("POST", "/academic-units", {
        "code": "VT-" + uuid.uuid4().hex[:8], "name": "Synthetic rejected unit"}, csrf=False)
    assert status == 403


def test_vt07_required_fields_rejected(admin):
    status, _, _ = admin.call("POST", "/careers", {})
    assert status == 400


def test_vt04_foreign_fixture_hidden():
    foreign_id = os.environ.get("SDD_FOREIGN_CAREER_ID")
    if not foreign_id:
        pytest.skip("A foreign synthetic fixture UUID is required")
    uuid.UUID(foreign_id)
    reader = Client()
    reader.login("SDD_READER")
    try:
        status, _, _ = reader.call("GET", "/careers/" + foreign_id)
        assert status == 404, "Fixture reader must not hold catalog admin or scope to target career"
    finally:
        reader.call("DELETE", "/sessions")


def test_vt12_logout_has_no_body_and_revokes():
    client = Client()
    client.login("SDD_ADMIN")
    status, _, body = client.call("DELETE", "/sessions")
    assert status == 204 and body is None
    assert client.call("GET", "/careers")[0] == 401

\end{Verbatim}
\section{Declaración de uso de Inteligencia
Artificial}\label{APE4-semana-4-declaraciuxf3n-de-uso-de-inteligencia-artificial}

ChatGPT/Codex apoyó revisión de fuentes, consulta oficial, elaboración
de pruebas/documentos y comprobaciones estáticas de artefactos. No se
ejecutó código del target ni se registraron feedback, commits, pruebas
satisfactorias o cambios funcionales inexistentes. La revisión personal
de los integrantes sigue pendiente.

\end{document}
